\chapter*{Résumé}
\addcontentsline{toc}{chapter}{Résumé}

Ce rapport présente le dimensionnement, la planification et l'optimisation d'un réseau LTE (bande~3, \SI{1800}{\mega\hertz}, porteuse de \SI{20}{\mega\hertz}) sur une zone réelle de \num{10} × \SI{10}{\kilo\metre} autour de Zaghouan (Tunisie), qui réunit un centre urbain dense, une campagne et le massif du Djebel Zaghouan. Dans une première étape, un dimensionnement analytique fixe les marges (effet de masque log-normal, pénétration, interférence), établit le bilan de liaison avec le modèle COST-231 Hata et en déduit le rayon des cellules : \SI{0,64}{\kilo\metre} en ville (indoor) et \SI{2,06}{\kilo\metre} en campagne (outdoor), la liaison montante étant limitante. La couverture impose 21 sites tri-sectoriels, alors que la capacité n'en demanderait que 5 dans le centre ; la réutilisation de fréquences retenue est $K = 1$.

Dans une seconde étape, ce réseau est construit dans Atoll 3.3 sur des cartes réelles (relief Copernicus, occupation du sol ESA WorldCover, population WorldPop) avec le modèle Cost-Hata et la diffraction sur le relief. Les prédictions de niveau de signal, de serveur, de recouvrement, de C/(I+N) et de débit, puis des simulations Monte-Carlo à 10, 100 et \num{1000} utilisateurs, montrent un réseau couvert en ville (\SI{99,5}{\percent} en indoor) mais troué derrière le massif (\SI{83,4}{\percent} de la campagne en outdoor), et dont les rejets d'utilisateurs viennent du signal et non de la capacité. Huit actions d'optimisation, mesurées une à une (tilt, retrait d'un site, deux nouveaux sites, ACP, ANR), portent la couverture rurale à \SI{94,3}{\percent} de la surface (\SI{97,9}{\percent} de la population) et la couverture urbaine à \SI{99,7}{\percent}, avec 21 sites, le nombre prévu par le calcul. Sur les mêmes \num{1041} utilisateurs simulés, les rejets passent de 108 à 71. Le calcul analytique a donc donné le bon nombre de sites, et la simulation sur le relief réel leurs bons emplacements.

\medskip
\textbf{Mots-clés :} LTE, planification radio, bilan de liaison, COST-231 Hata, Atoll, Monte-Carlo, DCFI, optimisation, ACP.

\vspace{2em}
\section*{Abstract}

This report covers the dimensioning, planning and optimisation of an LTE network (band~3, 1800~MHz, 20~MHz carrier) over a real 10 × 10~km area around Zaghouan, Tunisia. An analytical link budget with the COST-231 Hata model gives uplink-limited cell radii of 0.64~km (urban, indoor) and 2.06~km (rural, outdoor), hence 21 tri-sector sites with frequency reuse 1. The network is then built in Atoll 3.3 on real terrain, land-cover and population data. Coverage, interference and throughput predictions, followed by Monte-Carlo simulations with 10, 100 and 1000 users, show that the remaining problem is coverage behind the mountain, not capacity. Eight measured optimisation actions raise rural outdoor coverage from 83.4\,\% to 94.3\,\% while keeping 21 sites, the number predicted by the analytical step.
